\documentclass[letterpaper,10pt,conference]{ieeeconf}
\IEEEoverridecommandlockouts
\overrideIEEEmargins
\usepackage{amsmath,amssymb,booktabs,graphicx,url}
\title{Motion2Scene: When Executed-Motion Contrasts Do Not Improve Humanoid Decisions}
\author{Anonymous submission}
\begin{document}
\maketitle
\thispagestyle{empty}\pagestyle{empty}
\begin{abstract}
Motion-conditioned scene construction can concentrate physical labeling on
encounters that require adaptation without teaching a policy to select between
adaptations. We examine this distinction for humanoid overhead traversal using
a frozen tracker, seven qualified schedules, and a causal three-phase selector.
In a completed five-arm development comparison, executed-contrast construction
and feasibility-screened uniform each pass five of six contexts at each of
three acquisition seeds. Thirteen of fifteen learned policies choose a constant
schedule across these contexts. Historical replay changes the fitted parameters
but not the executed decisions. A retrospective screen audit recovers two of
seven measured response-complementarity positives; its two selected positives
do not establish population precision. These comparisons are consistent with
positive-support filtering explaining part of the construction advantage over
unscreened placement, but do not isolate a sole causal mechanism. A separately
registered support-broadening pilot has completed acquisition; its validation
panel remains incomplete after an interrupted execution. We report that status
without drawing a pilot utility conclusion. The evidence motivates separating
proposal geometry, physical labels, and downstream decision utility, while
establishing neither population equivalence nor held-out or hardware benefit.
\end{abstract}

\section{Introduction}
A humanoid that tracks walking and crouching motions still needs to decide when
and how to adapt beneath an obstacle. Motion-conditioned scene construction
addresses the collection problem: place constraints where an executed response
is predicted to clear and another is predicted to interfere. This can improve
the yield of physical labels. The consequential uncertainty is whether those
labels teach better decisions than a simpler constructor that preserves
predicted feasibility and samples uniformly.

The distinction matters because individually useful contrasts can be redundant
as a corpus. If one sustained response solves every training encounter, a
constant policy can fit the preferred response without using its observations.
Conversely, a policy can choose several schedules and still choose the wrong
one where their physical capabilities differ. Response diversity is therefore
a diagnostic of contextual selection, not a substitute for passage utility.
A constant policy could also improve passage; non-constancy is not necessary
for every possible utility gain.

We investigate these possibilities in a fixed humanoid simulation interface.
The completed development comparison does not show incremental passage benefit
over feasibility-screened uniform. This finite result revises the contribution
from a general data-utility claim to an analysis of the gap between geometric
contrast, physical response complementarity, and learned selection. We retain
the strongest script, all forced schedules, and reverse cases. The new pilot
asks whether broader proposal support changes this result, but its incomplete
outcomes cannot yet answer that question.

\section{Related Work}
Learning from Hallucination (LfH) constructs constrained navigation examples
from motion in open environments~\cite{lfh}. Learning from Learned Hallucination
(LfLH) learns obstacle configurations conditioned on existing plans and trains
reactive navigation policies~\cite{lflh}. These works motivate motion-conditioned
proposals. Dyna-LfLH extends learned hallucination to dynamic obstacles for
ground navigation~\cite{dyna}. Their navigation results do not validate this
humanoid pipeline or transfer their safety properties to it.

Our distinction is experimental: executed humanoid alternatives propose
constraints, obstacle-present closed-loop executions supply labels, and a fixed
selector tests the resulting data. An empty-space envelope does not certify
obstacle-present motion. A differentiable planner's reconstruction objective
cannot replace the physical passage labels used here. The low-level tracker
remains SONIC~\cite{sonic}; this study tests construction and selection within
its qualified command interface, not a newly trained locomotion controller.

\section{Decision Interface and Physical Outcome}
Let $q$ denote a scene, $s$ a qualified complete schedule, and $K$ the frozen
tracker. The measured response is
\begin{equation}
 \tau_s=\operatorname{Execute}(K,x_0,h_0,s,q,\xi),
\end{equation}
where $x_0,h_0$ specify the initial state and command history and $\xi$ is the
simulation seed. The seven schedules comprise neutral walking, early and late
short crouches, early and late sustained crouches, and two entries of a
projected/spliced prior-derived reference. Authored crouches are not independently
validated crawling skills. The interface permits one adaptation followed by
recovery on a fixed route, with no stop or steering action.

\begin{table}[t]\centering\small
\begin{tabular}{lcc}\toprule
Schedule & Entry (s) & Return (s)\\\midrule
Neutral & --- & ---\\
Short, early / late & 0.30 / 1.40 & 5.30\\
Sustained, early / late & 0.30 / 1.40 & 5.10\\
Prior derivative, early / late & 0.30 / 1.00 & 5.30\\\bottomrule
\end{tabular}
\caption{The complete seven-schedule bank. Each paired entry is a separate
qualified schedule. These times define the command interface, not guaranteed
obstacle clearance.}\label{tab:bank}
\end{table}

At reference ticks $15,50,70$ (0.30, 1.00, 1.40 s), a causal 114-dimensional
floor/ceiling-history and robot-state input supports a choice of a legal schedule
or WAIT. After commitment, the schedule specifies return. A WAIT target must
include a complete future continuation rather than crediting an unfinished
prefix. All construction arms share the complete-continuation teacher and the
phase ridge learner with $\lambda=10$. The five-arm comparison consequently
does not isolate the teacher's contribution.

The 299-sample reference permits 298 recorded control rows at 50 Hz and 1,192
physics substeps at 200 Hz. Contact evidence is retained at the physics rate.
The registered criterion requires initial upstream approach, ordered crossing
by all recorded body origins beyond the downstream face plus 0.1 m, and at least 0.30 s of upright stabilization.
Upright means root height at least 0.5 m and negative projected-gravity $z$
at least 0.5. Crossing is an origin-based criterion, not a certificate for every
native collision shape. Beam contact uses the maximum normal-force norm across
recorded bodies and time; the limit is 1.0 N. Multiple beams and other disallowed
environment contacts are audited explicitly. Permitted foot--floor support and
self-contact are recorded separately.

A task pass additionally requires the registered recovery and acceptable
contact and stability through the captured horizon. Insufficient stabilization
before that horizon is a task failure. A verified early physical failure is
retained even if later teaching targets are unavailable. Missing or unverifiable
measurements remain unknown. An operating-system exit code alone does not
classify physical success or technical failure. Successful passage time ends at
crossing plus stabilization; it is distinct from whole-episode duration and
recovery completion. No measured actuator work or energy claim is available.

\section{Construction and the Completed Development Comparison}
Executed contrasts use achieved empty-space body envelopes to rank candidate
constraints with both a positive robustness gate and a negative interference
gate. The positive test samples 81 geometric offsets. The negative test uses
the declared nominal interference margin. These tests are proposal heuristics:
obstacles may change tracking, contact, and the subsequent trajectory. Each
selected encounter receives the same student execution and seven complete
teacher branches. Outcomes, including all-fail response tables, remain in the
acquisition record.

The Gen~2 comparison evaluates five construction/training arms on six shared
development contexts at three acquisition seeds. The physics seed is common.
Screened Uniform retains predicted feasibility; Target-only is the weaker
unscreened placement comparator in this study. The other arms use reference
contrast, executed contrast, and executed contrast with historical replay.
These arms must not be conflated with the strict, broad, and mixture arms of
the subsequent pilot.

\begin{table}[t]\centering\small
\footnotesize
\begin{tabular}{lccc}\toprule
Learned arm at M8 & Seed 1 & Seed 2 & Seed 3\\\midrule
Screened Uniform & 5/6 & 5/6 & 5/6\\
Target-only & 4/6 & 4/6 & 5/6\\
Reference contrast & 6/6 & 5/6 & 4/6\\
Executed contrast & 5/6 & 5/6 & 5/6\\
Executed contrast + replay & 5/6 & 5/6 & 5/6\\\midrule
Strong script & \multicolumn{3}{c}{6/6}\\
Neutral & \multicolumn{3}{c}{2/6}\\
Short, early; late & \multicolumn{3}{c}{3/6; 3/6}\\
Sustained, early; late & \multicolumn{3}{c}{4/6; 4/6}\\
Prior derivative, early; late & \multicolumn{3}{c}{5/6; 5/6}\\
Union of seven fixed schedules & \multicolumn{3}{c}{6/6}\\\bottomrule
\end{tabular}
\caption{Completed six-context development panel: 138 executions, including
90 learned-policy, six script, and 42 forced-schedule assignments. Seed columns
index acquisition seeds 93201--93203; all executions use seed 8732. Contexts are
shared, so the cells are not independent experimental replicates.}
\label{tab:development}
\end{table}

Executed contrast and Screened Uniform pass the same five contexts at every
acquisition seed. Their paired passage difference is zero, with zero wins,
zero losses, and eighteen ties across matched seed--context cells. Passage-time
differences on mutually successful cells are also zero. This is no observed
incremental benefit on this finite development panel; it is neither population
equivalence nor evidence of non-inferiority. There is no registered
non-inferiority margin or suitable uncertainty bound established here.

The development panel was used during method development. In particular, its
late complementary context was generated using a clearance predicate and
ranking related to the constructor. It cannot provide independent confirmation
of that constructor. We report the entire panel rather than a favorable subset.
The strong script passes all six contexts, and no learned policy exceeds the
measured capability of the seven-schedule bank.

\section{Mechanism Diagnostics and Counterexamples}
Thirteen of fifteen learned policies select a single schedule across the six
contexts. The two non-constant policies arise from corpora requiring more than
one schedule for coverage. This correspondence was noticed after inspecting
outcomes; shared seeds and contexts also invalidate treating the fifteen
corpora as independent observations for a simple association test. It is a
post-hoc association, not proof that complementarity causes contextual learning.

The reverse cases constrain the explanation. The third Target-only policy uses
three schedules but still scores 5/6, selecting a failing short response on the
late complementary context. The first reference-contrast policy reaches 6/6;
its final acquired encounter was rejected by the executed positive robustness
bar. Thus neither non-constancy nor passing the executed geometric gates is a
necessary-and-sufficient description of useful decision data.

Across 92 already executed encounters, the deployed negative screen identifies
two of seven measured positives for the property that covering schedules fail
while another schedule passes. Both selected cases are positive. The resulting
2/2 precision is retrospective and selection-conditioned, not an estimate of
established population precision. Relaxing the negative margin from $-10$ mm to
zero recovers only one additional measured positive. Some missed encounters
have predicted positive clearance for responses that physically fail; another
is excluded by the positive gate. Negative-threshold relaxation therefore does
not resolve every observed omission.

Historical replay is a measured null on this panel. It changes fitted weights
and biases while preserving every chosen schedule and outcome in the eighteen
matched executions. Constancy here does not imply that the ridge model is
structurally incapable of conditioning on observations. Complete-continuation
teaching was common to the comparison and cannot receive causal credit from it.
The positive-support explanation is consistent with the observed advantage
over the weaker placement comparator and the tie with Screened Uniform. It
remains one explanation among coupled changes in acquired data.

\section{Registered Support-Broadening Pilot: Incomplete}
The pilot preserves the original three proposal arms:
\begin{itemize}
\item A, strict: both unchanged executed-contrast gates and ranking.
\item B, broad: valid task-domain candidates in draw order, with neither
geometric gate applied.
\item C, mixture: the committed strict/broad/strict/broad alternation.
\end{itemize}
The strict reservoir admission rate is approximately 6.4\%. B does not preserve
the positive screen. C preserves broad proposal support, not universal geometric
certification. All arms share, within each acquisition seed, the byte-identical
four-encounter prefix and M4 checkpoint, frozen tracker, seven schedules, three
phases, causal inputs, teacher, learner, and physics settings. The assignment
ledger declares 288 new acquisition episodes over nine corpora, with four new
encounters and eight executions per encounter. These totals are assignment
budgets, not measured execution cost.

Ten validation contexts were drawn by beam family and underside band using
task geometry. The generation manifest reports zero clearance queries and zero
physics steps and excludes physical outcomes from generation. Contexts are
retained regardless of solvability and kept outside training and proposal
selection. Eighteen reserved layouts remain outside this pilot. A contemporary
hash verifies artifact integrity but cannot prove when registration occurred.

The registered validation panel has 200 assignments: nine pilot policies,
three reused M8 executed-contrast policies, the strong script, and seven fixed
schedules on all ten contexts. It omits a contemporaneous Screened Uniform
policy. The six-context 5/6 score cannot fill that missing comparison. A baseline
utility test would require a separately registered and budgeted amendment;
no episodes are silently added here.

At the September 11 audit boundary, acquisition has 288 scored assignment
receipts. Validation has 171 scored receipts, an interrupted assignment with
no terminal attempt receipt or trajectory/contact streams, and 28 unrun
assignments. The frozen runner refuses automatic retry of the interrupted
launch. Execution is therefore blocked. We withhold pilot utility results
until complete reconciliation; the partial panel is not a substitute for the
registered comparison. No new physics was launched during this audit.

The registered contrasts are B versus A and C versus A for broadening, and C
versus B for guided-half value. Improvement over A cannot isolate removal of the
negative gate because B also removes the positive gate. A C--B tie would not
establish equivalence. Exact permutation inference would require an explicit
exchangeability scheme respecting shared contexts, matched seeds, and corpora.
No such scheme or multiplicity rule has been established for this incomplete
panel, so inferential $p$-values are NA. Ten sampled contexts and three acquisition
seeds cannot be reinterpreted as 200 independent replications. New audit
comparisons are labeled exploratory after access to partial outcomes.

\section{Scope, Adoption Gates, and Conclusion}
The completed development evidence establishes a finite null for incremental
passage over Screened Uniform. It leaves the broader population effect
inconclusive. The pilot cannot yet support either broadening or guided-half
utility, and its omitted baseline prevents a validation claim beyond
feasibility-screened placement. More diverse schedules alone would not change
these conclusions. A dataset utility claim additionally needs separate reuse,
licensing, and reproduction evidence.

Phase~3 remains closed. Before any reserved outcomes are opened, implementation
and terminal-censoring policies must be locked and all adoption gates satisfied.
Source-ancestry-held-out transfer must precede generalization claims. Realistic
sensing, noise and dropout, timing, and a proposed $70^\circ$ field of view require
a separate sensor specification and registration. The latter is not a measured
deployed configuration. Delays of 100, 260, and 500 ms require matched sensing
conditions; small historical panels do not define universal delay limits.
Sequential decisions need matched aggregation and fresh labels across arms,
qualification of every skill and transition, and whole-chain success. The
existing crouch bank does not establish crawling, recovery chains, or steering.

All physical evidence here is simulation. A MuJoCo visualization of a recorded
trajectory is not a second dynamics experiment. Refusal in the recorded interface
continues walking, and the 1 N contact threshold is a simulation convention.
Neither supports hardware safety or protective stopping. Motion-conditioned
construction is useful as a proposal mechanism; its downstream value must be
measured against a strong screened comparator rather than inferred from
geometric contrast or training-label diversity.

\section*{Acknowledgments}
OpenAI Codex (GPT-6) assisted with the September 11 evidence audit and generated
revision text throughout this manuscript, including its abstract, interpretation,
and limitations. It also generated the local reconciliation code. Submission
requires author verification of the evidence, citations, and final text.

\begin{thebibliography}{9}
\bibitem{lfh} X. Xiao, B. Liu, G. Warnell, and P. Stone,
``Toward agile maneuvers in highly constrained spaces: Learning from hallucination,''
\emph{IEEE Robotics and Automation Letters}, vol. 6, no. 2, pp. 1503--1510, 2021.
\bibitem{lflh} Z. Wang et al., ``From agile ground to aerial navigation:
Learning from learned hallucination,'' in \emph{IROS}, pp. 148--153, 2021.
\bibitem{dyna} S. A. Ghani, Z. Wang, P. Stone, and X. Xiao,
``Dyna-LfLH: Learning agile navigation in dynamic environments from learned
hallucination,'' arXiv:2403.17231, 2024.
\bibitem{sonic} Z. Luo et al., ``SONIC: Supersizing motion tracking for natural
humanoid whole-body control,'' arXiv:2511.07820, 2025.
\end{thebibliography}
\end{document}
